\documentclass[10pt]{article}
\usepackage[utf8]{inputenc}
\usepackage[T1]{fontenc}
\usepackage[margin=0.5in]{geometry}
\usepackage{amsmath, amssymb, bm}
\usepackage[colorlinks=true,citecolor=blue,linkcolor=blue]{hyperref}
\usepackage{enumitem}
\usepackage{setspace}

\setlength{\parskip}{6pt}
\setlength{\parindent}{0pt}
\setstretch{0.1}
\renewcommand{\refname}{}

\title{\vspace{-1.5em}\textbf{Speculative Inference for Masked Diffusion Language Models}\vspace{-1em}}
\author{Pranav Sitaraman (pranavsi@mit.edu) and Gavin Ye (gavinye@mit.edu)}
\date{}

\begin{document}
\maketitle
\thispagestyle{empty}
\vspace{-1.5em}
\textbf{Team.} As we already discussed with course instructors both in person and over email, Gavin and I have been working for a couple weeks on an NLP-related research project with help from a Harvard grad student friend of ours, aiming to improve inference-time steering and KV caching for masked diffusion LLMs. Although our third friend is not taking the class, we will bear all responsibility for the success/outcome of the project, write the entirety of the final report and poster presentation ourselves, and take credit for only the parts of the research that we directly work on. We will also present between the two of us and will work on most of the theoretical design and experimentation ourselves. The project is entirely self-contained and only relies on outside mentorship to guide the literature review and brainstorming processes.

\textbf{Direction.} Masked diffusion language models (dLLMs) generate text by iteratively unmasking tokens in parallel, offering any-order properties that autoregressive models lack. However, frontier dLLMs currently adopt inference strategies like fixed confidence thresholds or rigid block schedules that fail to exploit this flexibility, with significant speed and quality gains still possible. We propose a method for \textbf{Any-Order Adaptive Editing}, a framework that unifies three capabilities that have so far been studied only in isolation: (i) learned adaptive unmasking, (ii) self-correction via remasking, and (iii) speculative KV-caching. The core idea is to pair a fast hard-routed auxiliary model with a slow soft-routed primary, both instantiated from the same frozen Mixture-of-Experts (MoE) dLLM, in a speculative diffusion loop where the auxiliary drafts tokens and pre-caches KV states, the primary verifies and refines, and a policy decides per-position unmask, remask, and cache actions.

\textbf{Motivation.} Three primary findings motivate this work. First, learned unmasking policies trained via reinforcement learning outperform heuristic samplers on diffusion LLMs but address only where to unmask, not self-correction or caching. Second, PRISM demonstrates that remasking low-quality tokens and regenerating them improves output quality but relies on a fixed threshold rather than a learned policy. Third, diffusion-aware KV-caching methods (dKV-Cache, Elastic-Cache) achieve substantial speedups but constrain the decoding trajectory to suit the cache, sacrificing the any-order property. Our paradigm is designed to combine all three under a single policy trained end-to-end with GRPO, while a composed prediction mechanism biases token selection toward cache-aligned choices to maximize KV reuse.

\textbf{Tentative Plan.}
\begin{enumerate}[leftmargin=*,topsep=0pt]
    \item Integrate the dual-model MoE setup (hard-routed auxiliary and soft-routed primary both from LLaDA 2.1-mini) into the dInfer/SGLang serving stack. Verify that soft routing with varying temperature $\tau$ produces the expected agreement tradeoff between auxiliary and primary models.
    \item Train the PRISM quality adapter. Implement soft-masked state construction, the lightweight transformer policy with three heads (unmask, remask, cache), and composed prediction and speculative KV-caching.
    \item Train via GRPO with a correctness and cache use/thrashing reward. Validate on GSM8K before scaling to MATH and HumanEval. Sweep the $\tau$ Pareto frontier of inference speed even as auxiliary/primary model agreement degrades.
    \item Run architecture ablations (removing auxiliary, policy, composed prediction, remask head, agreement signal). Evaluate positional variants (learned access-pattern prediction, layer-wise refresh depth, history-conditioned features) and attention alternatives for the auxiliary model (sliding-window, sparse, linear). Compare full speculation (caching location and token) against location-only steering with training-free token selection methods.
\end{enumerate}

\textbf{Anticipated Risks and Challenges.}

\begin{itemize}[leftmargin=*, topsep=0pt]
    \item We rely on the premise that using all 16B parameters via soft routing yields meaningfully better predictions than the 1.4B hard-routed subset. We can also consider fine-tuning the soft-routed variant to ensure the dual-model architecture's benchmark performance compares to single-model baselines. We will test this early before committing to the full pipeline.
    \item Prior work trained a single Bernoulli head (unmask only). Adding remask and cache heads triples the action dimension and introduces conflicting objectives (speed reward favors aggressive caching, while the thrash penalty punishes caching positions that are later remasked). The policy may converge to degenerate strategies, so careful reward shaping and staged training (unmask-only first, then adding heads) may be needed. This challenge of credit-assignment may also inevitably degrade model performance during GRPO.
\end{itemize}

\textbf{References.}

\bibliographystyle{plain}
\begin{thebibliography}{99}
\vspace{-27pt}

\bibitem{nie2025llada2}
S.~Nie et al. Large language diffusion models. \emph{arXiv:2502.09992}, 2025.

\bibitem{llada21}
S.~Nie et al. {LLaDA 2.1}: Speeding up text diffusion via token editing. \emph{arXiv:2602.08676}, 2025.

\bibitem{jazbec2025learning}
M.~Jazbec et al. Learning unmasking policies for diffusion language models. \emph{arXiv:2512.09106}, 2025.

\bibitem{kim2025trainworst}
J.~Kim et al. Train for the worst, plan for the best: Understanding token ordering in masked diffusions. \emph{ICML}, 2025.

\bibitem{kim2025prism}
J.~Kim et al. {PRISM}: Plug-in remasking for inference-time self-correction of masked diffusions. \emph{arXiv:2510.01384}, 2025.

\bibitem{ma2025dkv}
Y.~Ma et al. {dKV-Cache}: The cache for diffusion language models. \emph{NeurIPS}, 2025.

\bibitem{elastic_cache}
Q.~Nguyen-Tri et al. Attention is all you need for KV cache in diffusion LLMs. \emph{ICLR}, 2026.

\bibitem{hersche2025softmasked}
M.~Hersche et al. Soft-masked diffusion language models. \emph{arXiv:2510.17206}, 2025.
\end{thebibliography}
\end{document}
